\documentclass[11pt]{article}

\usepackage[margin=1in]{geometry}
\usepackage{amsmath,amssymb,amsfonts}
\usepackage{amsthm}
\usepackage{algorithm}
\usepackage{algorithmic}
\usepackage{booktabs}
\usepackage{hyperref}
\usepackage{graphicx}
\usepackage{enumitem}
\usepackage{xcolor}

\hypersetup{colorlinks=true, linkcolor=blue!45!black, citecolor=blue!45!black,
  urlcolor=blue!45!black}

\newtheorem{assumption}{Assumption}
\newtheorem{remark}{Remark}
\newtheorem{definition}{Definition}

\newcommand{\Dold}{\mathcal{D}_{\mathrm{old}}}
\newcommand{\Dnew}{\mathcal{D}_{\mathrm{new}}}
\newcommand{\X}{\mathbf{X}}
\newcommand{\z}{\mathbf{z}}
\newcommand{\POrisk}{\mathrm{PO\text{-}risk}}
\newcommand{\Rrisk}{\mathrm{R\text{-}risk}}
\newcommand{\VIMP}{\mathrm{VIMP}}
\newcommand{\MMD}{\mathrm{MMD}}
\newcommand{\HSIC}{\mathrm{HSIC}}
\newcommand{\CS}{\textsc{cs}}
\newcommand{\CD}{\textsc{cd}}
\newcommand{\Exp}{\mathbb{E}}

\title{\bf Feature Selection for Distribution Shift in Structured Reasoning:\\
Monitoring and Steering Skeleton-of-Thought and Tree-of-Thoughts}
\author{}
\date{}

\begin{document}
\maketitle

\begin{abstract}
Structured reasoning frameworks such as Skeleton-of-Thought (SoT) and
Tree-of-Thoughts (ToT) substantially improve the quality of large-language-model
reasoning, but their behaviour is acutely sensitive to task distribution, prompt
template, model version, and domain. When the workload drifts, these systems
degrade silently: accuracy, latency, and cost move without an obvious cause, and
existing observability tools treat a reasoning trace as an opaque token stream
that offers no purchase on \emph{where} or \emph{why} the drift occurred. We
adapt Feature Selection for Distribution Shift (FSDS)---an assumption-less,
data-generating-process-agnostic procedure originally developed for tabular and
multi-modal feature attribution---to structured reasoning. We represent a trace
as a graph- or tree-structured feature object with node-, edge-, and
trace-level covariates, and cast the comparison of a reference batch against a
live batch as a two-sample problem in which the batch label plays the role of a
treatment indicator. FSDS then distinguishes covariate shift in $P(\X)$ from
concept drift in $P(Y\mid\X)$, ranks the features that drive the shift through
out-of-bag variable importance, MMD/HSIC leave-one-covariate-out, and
meta-learner risk (CFPerm, PermuCATE, DRPerm, and PO-risk LOCO), and converts
the top-ranked features into concrete interventions: re-plan the skeleton,
re-expand a node, adjust the branching or pruning budget, re-calibrate the value
model, route to a stronger model, or fall back to a cheaper policy. The
procedure is model-agnostic, robust to missing modalities and missing reasoning
branches, and equipped with an automatic fallback under poor overlap, making it
suitable for continuous production monitoring. We emphasise throughout that the
meta-learners are used strictly for distance estimation, not for causal
identification.
\end{abstract}

\section{Introduction}
\label{sec:intro}

Chain-of-Thought prompting \cite{wei2022cot} established that eliciting
intermediate reasoning steps improves the performance of large language models
(LLMs) on multi-step tasks. Subsequent frameworks structure this computation
explicitly: Skeleton-of-Thought (SoT) \cite{ning2024sot} first generates a
skeleton of reasoning points and then expands them, frequently in parallel, to
reduce latency; Tree-of-Thoughts (ToT) \cite{yao2023tot} searches over a tree of
intermediate thoughts with expansion, evaluation, pruning, and backtracking to
improve deliberate problem solving. These paradigms are effective, but their
behaviour is highly non-stationary in deployment. The distribution of incoming
tasks changes, prompt templates are revised, the underlying model is upgraded,
and the retrieval corpus evolves. Any of these can shift both the inputs the
reasoner sees and the mapping from a reasoning trace to a correct answer.

The resulting failures are difficult to diagnose. Aggregate metrics such as
end-to-end accuracy or mean latency reveal \emph{that} something changed, but not
\emph{which} part of the reasoning process is responsible, nor whether the change
is a shift in the workload (which may warrant re-planning or re-routing) or a
degradation in the mapping from reasoning to outcome (which may warrant
re-calibration or re-expansion). Treating the trace as an opaque sequence
forecloses this distinction.

We argue that a reasoning trace is naturally a \emph{structured feature object}
and that the machinery of feature-level distribution-shift attribution applies
directly to it. Feature Selection for Distribution Shift (FSDS) was developed to
identify, without assumptions on the data-generating process, the subset of
features that drive a shift between two batches of data, and to separate
covariate shift ($P(\X)$) from concept drift ($P(Y\mid\X)$)
\cite{fsds2026}. Its estimators---out-of-bag (OOB) variable importance from a
random-forest domain classifier, maximum-mean-discrepancy and
Hilbert--Schmidt-independence-criterion leave-one-covariate-out (MMD/HSIC-LOCO),
and meta-learner risk functionals (CFPerm, PermuCATE, DRPerm, PO-risk)---are
model-agnostic, compatible with arbitrary modalities, robust across model
complexities, and cheap to run as $O(p)$ post-fit operations. These are exactly
the properties required to monitor a reasoning system whose feature space mixes
embeddings, discrete search statistics, and cost counters.

\paragraph{Contributions.} This manuscript drafts a prototype and its evaluation
protocol. Concretely, we:
\begin{enumerate}[leftmargin=*,itemsep=2pt]
  \item formulate SoT and ToT reasoning traces as structured feature objects and
  give a node/edge/trace feature taxonomy suitable for FSDS
  (Sections~\ref{sec:formulation}--\ref{sec:features});
  \item specialise the FSDS attribution pipeline to the two frameworks, yielding
  \textsc{FSDS-SoT} and \textsc{FSDS-ToT}, which localise drift to skeleton
  positions, tree depths, branches, and value bins, and separate covariate shift
  from concept drift at each level (Sections~\ref{sec:pipeline}--\ref{sec:tot});
  \item define an attribution-to-intervention policy that maps the dominant drift
  type at each localised feature group to a concrete control action, with an
  automatic fallback under poor overlap (Section~\ref{sec:intervention}); and
  \item describe a reproducible evaluation protocol and a production deployment
  design, including chronological splitting, parallelised permutation, and a
  gated mixture over reasoning experts
  (Sections~\ref{sec:protocol}--\ref{sec:deployment}).
\end{enumerate}
We are explicit about scope: the meta-learners are used only as
distance/deviation estimators, and the placeholder attribution tables are
illustrative of the output format rather than empirical results.

\section{Related Work}
\label{sec:related}

\paragraph{Structured reasoning.} Beyond CoT \cite{wei2022cot}, SoT
\cite{ning2024sot} and ToT \cite{yao2023tot}, a family of frameworks structures
the reasoning computation in different ways: Graph-of-Thoughts
\cite{besta2024got} generalises the tree to a directed acyclic graph with
aggregation; ReAct \cite{yao2023react} interleaves reasoning with tool actions
and environment feedback; Self-Consistency \cite{wang2023selfconsistency}
marginalises over independently sampled chains; and Reflexion
\cite{shinn2023reflexion} adds iterative self-critique with verbal memory. All of
these produce segmentable intermediate representations, which is precisely what
our feature map exploits. Our contribution is orthogonal to the choice of
framework: we monitor and steer it rather than replace it.

\paragraph{Distribution-shift detection.} Two-sample testing via MMD
\cite{gretton2012mmd} and dependence testing via HSIC \cite{gretton2005hsic} are
standard tools for detecting covariate shift; model-X knockoffs
\cite{candes2018knockoffs} provide controlled variable selection. For concept
drift, meta-learners rooted in causal inference---the R-learner
\cite{nie2021rlearner}, the doubly-robust DR-learner \cite{kennedy2023drlearner},
and causal forests \cite{wager2018causalforest}---estimate a deviation between
conditional distributions. FSDS \cite{fsds2026} combines a random-forest domain
classifier \cite{breiman2001rf} with these functionals, using them as distance
estimators equipped with a permute-then-refit procedure, and demonstrates that
parsimonious random-forest importance is competitive with, and often more robust
than, more elaborate alternatives. We inherit these estimators wholesale and
apply them to reasoning-trace features.

\paragraph{LLM observability.} Production monitoring of LLM systems typically
tracks scalar quality and cost signals or uses univariate drift indices such as
the population stability index (PSI) or KL divergence per feature. These
approaches neither localise drift to a component of the reasoning process nor
separate input shift from mapping degradation. Our prototype provides both.

\section{Problem Formulation}
\label{sec:formulation}

\paragraph{Traces.} A SoT trace is an ordered set of skeleton points
$\mathcal{S}=\{s_1,\dots,s_m\}$, where each $s_j$ is expanded into a thought
$t_j$ and the trace yields a final answer $Y$. A ToT trace is a tree
$\mathcal{T}=(\mathcal{V},\mathcal{E})$ in which each node $v$ carries a state
embedding $\z_v$, a value $V(v)$, a depth $d(v)$, a branching factor $b(v)$, and
a prune indicator, and the answer is obtained by voting or best-path selection.

\paragraph{Feature map.} A trace is summarised by a feature map
\begin{equation}
  \X \;=\; \phi(\mathcal{T}\ \text{or}\ \mathcal{S}) \;\in\; \mathbb{R}^p,
  \label{eq:featuremap}
\end{equation}
whose coordinates are grouped by structural level (Section~\ref{sec:features}).
The outcome $Y$ is a task-appropriate quality signal: exact-match correctness, a
verifier or process-reward score, a unit-test pass rate, or a business KPI.

\paragraph{Batch-as-treatment.} Given a reference batch and a live batch,
\begin{equation}
  \Dold=\{(\X_i^{\mathrm{old}},Y_i^{\mathrm{old}})\}_{i=1}^{n_0}, \qquad
  \Dnew=\{(\X_i^{\mathrm{new}},Y_i^{\mathrm{new}})\}_{i=1}^{n_1},
\end{equation}
we assign a binary batch label $T\in\{0,1\}$ (old $=0$, new $=1$) and treat it as
a treatment indicator. FSDS then tests, and attributes at the feature level, the
two hypotheses
\begin{equation}
  \CS:\ P(\X_{\mathrm{new}})\neq P(\X_{\mathrm{old}}),
  \qquad
  \CD:\ P(Y_{\mathrm{new}}\mid\X_{\mathrm{new}})\neq P(Y_{\mathrm{old}}\mid\X_{\mathrm{old}}).
\end{equation}

\begin{assumption}[Overlap]
\label{ass:overlap}
The propensity $e(\X)=P(T=1\mid\X)$ is bounded away from $0$ and $1$ on the
support of interest, so that the two batches share common support. When this
fails, the concept-drift estimators are unreliable and the procedure falls back
as described in Section~\ref{sec:intervention}.
\end{assumption}

\begin{remark}[Distance estimation, not causal identification]
\label{rem:distance}
Although the concept-drift estimators borrow the R-learner, DR-learner, and
causal-forest machinery, we use them solely to estimate the deviation between
conditional distributions across batches. No causal interpretation of the batch
label or of individual features is claimed or required
\cite{fsds2026,nie2021rlearner,kennedy2023drlearner}.
\end{remark}

\section{Reasoning Traces as Feature Objects}
\label{sec:features}

The feature map $\phi$ in \eqref{eq:featuremap} is organised into six structural
levels---trace, node, edge, text, quality, and cost---each populated with
framework-specific covariates (Table~\ref{tab:taxonomy}). Node embeddings enter
as dense blocks; discrete search statistics (depth, branching, prune flags) enter
as scalars; text is encoded by a frozen sentence or multi-modal encoder. Because
FSDS is agnostic to feature type and tolerates missing coordinates, traces with
absent branches (a pruned subtree) or absent modalities (a text-only query in a
multi-modal system) are handled without imputation heuristics: the missingness
pattern is itself a feature.

\begin{table}[t]
\centering
\caption{Reasoning feature taxonomy for FSDS. Node embeddings enter as grouped
dense blocks; the remaining coordinates are scalar or categorical.}
\label{tab:taxonomy}
\small
\begin{tabular}{lll}
\toprule
Level & SoT features & ToT features \\
\midrule
Trace & skeleton length, parallel width, total tokens & tree depth, total nodes, token cost \\
Node & node order, node embedding, expansion length & node value, depth, branching, prune flag \\
Edge & dependency edge, order edge & parent--child edge, backtrack edge \\
Text & query, skeleton text, thought text & node text, value text, final-answer text \\
Quality & self-consistency, retrieval hit & value entropy, vote entropy, prune confidence \\
Cost & latency, FLOPs, API calls & expansion count, evaluation count, search cost \\
\bottomrule
\end{tabular}
\end{table}

\section{The FSDS Attribution Pipeline}
\label{sec:pipeline}

FSDS attributes shift along two complementary planes.

\paragraph{Covariate plane (no outcome required).} A random-forest domain
classifier is fit to predict the batch label $T$ from $\X$; its OOB permutation
variable importance ranks the coordinates that distinguish the batches, and its
AUC summarises the overall magnitude of covariate shift. In parallel, a
statistical route quantifies each feature's contribution through a
leave-one-covariate-out drop in the maximum mean discrepancy,
\begin{equation}
  \Delta\MMD_j \;=\; \MMD^2(\X^{\mathrm{old}},\X^{\mathrm{new}})
  \;-\; \MMD^2(\X^{\mathrm{old}}_{-j},\X^{\mathrm{new}}_{-j}),
  \label{eq:mmdloco}
\end{equation}
with an RBF kernel ($\sigma=1$); a large positive $\Delta\MMD_j$ indicates that
feature $j$ drives the covariate shift, since removing it collapses the
discrepancy. HSIC-LOCO is defined analogously for dependence-based attribution.
The covariate plane requires no outcome and is therefore always available, even
online.

\paragraph{Concept plane (outcome required).} To attribute concept drift we test
whether $P(Y\mid\X)$ differs across batches using CFPerm/PermuCATE and fit a
doubly-robust learner with the batch label as treatment. With nuisances
$m(\X)=\Exp[Y\mid\X]$ and $e(\X)=P(T=1\mid\X)$, the doubly-robust pseudo-outcome
\begin{equation}
  \psi_i \;=\; \hat\tau(\X_i)
  \;+\; \frac{T_i-\hat e(\X_i)}{\hat e(\X_i)\,\bigl(1-\hat e(\X_i)\bigr)}\,
  \bigl(Y_i-\hat m(\X_i)\bigr),
  \qquad
  \POrisk \;=\; \frac{1}{n}\sum_{i=1}^{n}\bigl(\psi_i-\hat\tau(\X_i)\bigr)^2,
  \label{eq:porisk}
\end{equation}
is consistent as a distance functional if either $m$ or $e$ is correctly
specified. Feature-level concept-drift attribution follows from a
leave-one-covariate-out increase in this risk,
\begin{equation}
  \phi_j^{\mathrm{LOCO}} \;=\; \POrisk(\X_{-j}) - \POrisk(\X),
  \label{eq:poloco}
\end{equation}
where $\X_{-j}$ uses the conditional fit $X_j\sim X_{-j}$ (random forest for
continuous coordinates, (multinomial) logistic regression for categorical ones).
Because $\POrisk$ is invariant to covariate shift, \eqref{eq:poloco} isolates the
contribution of feature $j$ to the change in the conditional mapping. Nuisances
are cross-fit (or estimated out-of-bag) to avoid overfitting bias, and the
random forest is the recommended specification for both $m$ and $e$ for
robustness.

\paragraph{Aggregation and overlap.} Per-coordinate scores are aggregated into
interpretable drift shares by structural group---skeleton position, node type,
and text source for SoT; depth, branch, and value bin for ToT. Every share is
accompanied by a bootstrap confidence interval and a top-$k$ stability
proportion. The propensity distribution, the domain-classifier AUC, and the
effective sample size monitor Assumption~\ref{ass:overlap}; when overlap is poor
the concept plane is suppressed and the procedure degrades to the covariate
plane or to univariate monitoring (Section~\ref{sec:intervention}).

\section{FSDS-SoT}
\label{sec:sot}

For SoT we treat each skeleton node as a candidate feature group and localise
drift to skeleton positions and node types. Covariate shift concentrated on
skeleton length or node order signals that the \emph{shape} of the decomposition
has changed---typically a workload shift that warrants re-planning the
skeleton---whereas concept drift concentrated on a specific node embedding
signals that the mapping from that reasoning point to the outcome has degraded,
warranting re-expansion of the node or routing to a stronger model. The full
procedure is given in Algorithm~\ref{alg:fsds-sot}; its cost is dominated by a
single domain-classifier fit and the meta-learner nuisances, with the LOCO passes
reduced to $O(p)$ post-fit operations.

\begin{algorithm}[t]
\caption{\textsc{FSDS-SoT}}
\label{alg:fsds-sot}
\begin{algorithmic}[1]
\REQUIRE reference traces $\Dold$, live traces $\Dnew$, budget $k$
\ENSURE top-$k$ drift nodes with covariate/concept shares and actions
\STATE Extract skeleton features $\X$ and outcomes $Y$ via $\phi$
\STATE Fit RF domain classifier on $(\X,T)$; compute OOB $\VIMP$ and AUC
\FOR{each skeleton node $j$}
    \STATE Compute $\Delta\MMD_j$ by MMD-LOCO \eqref{eq:mmdloco}
\ENDFOR
\STATE Run CFPerm/PermuCATE to test the $P(Y\mid\X)$ shift
\STATE Fit the DR learner; compute $\POrisk$ and $\phi_j^{\mathrm{LOCO}}$ \eqref{eq:porisk}--\eqref{eq:poloco}
\STATE Aggregate covariate/concept shares by node position and node type
\IF{$\CS$ dominates at skeleton length or order}
    \STATE Re-plan the skeleton
\ENDIF
\IF{$\CD$ dominates at node $j$}
    \STATE Re-expand node $j$ or route to a stronger model
\ENDIF
\IF{overlap $<0.3$ or $n<50$}
    \STATE Fall back to RF-domain $\VIMP$ or univariate PSI/KL
\ENDIF
\RETURN top-$k$ nodes with shares and actions
\end{algorithmic}
\end{algorithm}

\section{FSDS-ToT}
\label{sec:tot}

For ToT we treat tree nodes and search statistics as features and localise drift
to depths, branches, and value bins (Algorithm~\ref{alg:fsds-tot}). Two patterns
are of particular operational interest. Concept drift concentrated on the value
signal indicates that the value model has become miscalibrated relative to the
new workload, warranting re-calibration of the value head rather than a change to
the search budget. Covariate shift concentrated on depth or branching indicates
that the search is exploring a differently shaped space, warranting an adjustment
of the branching factor, maximum depth, or prune threshold. As in the SoT case, a
poor-overlap fallback protects the procedure when the two batches are too easily
separable or the live batch is too small.

\begin{algorithm}[t]
\caption{\textsc{FSDS-ToT}}
\label{alg:fsds-tot}
\begin{algorithmic}[1]
\REQUIRE reference traces $\Dold$, live traces $\Dnew$, budget $k$
\ENSURE top-$k$ drift features with shares and a ToT policy
\STATE Extract node features: depth, value, branching, prune flag, path entropy, cost
\STATE Fit RF domain classifier on $(\X,T)$; compute OOB $\VIMP$
\STATE Run CFPerm/PermuCATE for the $P(Y\mid\X)$ shift
\STATE Fit the DR learner; compute $\POrisk$ and per-feature $\phi_j^{\mathrm{LOCO}}$
\STATE Aggregate covariate/concept shares by depth, branch, and value bin
\IF{$\CD$ dominates at the value signal}
    \STATE Re-calibrate the value head
\ENDIF
\IF{$\CS$ dominates at depth or branching}
    \STATE Adjust branching factor, maximum depth, and prune threshold
\ENDIF
\IF{overlap $<0.3$}
    \STATE Fall back to SoT/CoT or univariate monitoring
\ENDIF
\RETURN top-$k$ features with shares and a ToT policy
\end{algorithmic}
\end{algorithm}

\section{From Attribution to Intervention}
\label{sec:intervention}

The value of localised attribution is that it maps to a small, well-defined set
of control actions. Table~\ref{tab:drift} illustrates the format of the drift
attribution output---per feature group, the covariate and concept shares and the
recommended action---and Table~\ref{tab:policy} summarises the intervention
policy as a function of the signal and its dominant drift type. Both tables are
illustrative of the output schema; they are not empirical results.

The governing logic is a separation principle. When covariate shift dominates,
the workload has moved and the reasoning \emph{structure} should adapt: re-plan
the skeleton, adjust the search budget, or re-route. When concept drift
dominates, the mapping from reasoning to outcome has degraded and the reasoning
\emph{content or scoring} should adapt: re-expand a node, re-calibrate the value
model, or change the prune threshold. When Assumption~\ref{ass:overlap} is
violated---diagnosed by an overlap below $0.3$ or a live batch smaller than a
threshold---no confident attribution is possible, and the procedure abstains,
falling back to the covariate plane or to univariate monitoring and issuing a
conservative default.

\begin{table}[t]
\centering
\caption{Illustrative drift attribution output (schema, not results).}
\label{tab:drift}
\small
\begin{tabular}{lccc}
\toprule
Feature group & $\CS$ share & $\CD$ share & Action \\
\midrule
Skeleton length & 0.22 & 0.08 & re-plan skeleton \\
Node 3 embedding & 0.10 & 0.31 & re-expand node \\
Tree depth & 0.18 & 0.12 & adjust maximum depth \\
Branch value & 0.07 & 0.28 & re-calibrate value head \\
Text source & 0.15 & 0.19 & route / update prompt \\
\bottomrule
\end{tabular}
\end{table}

\begin{table}[t]
\centering
\caption{Illustrative intervention policy (schema, not results).}
\label{tab:policy}
\small
\begin{tabular}{lll}
\toprule
Signal & Dominant drift & Policy \\
\midrule
Skeleton order & $\CS$ & re-plan skeleton \\
Node value & $\CD$ & re-calibrate value model \\
Depth / branching & $\CS$ & adjust search budget \\
Prune flag & $\CD$ & change prune threshold \\
Missing branch & either & fall back to SoT/CoT \\
\bottomrule
\end{tabular}
\end{table}

\section{Experimental Protocol}
\label{sec:protocol}

We describe the protocol by which the prototype is to be validated; empirical
results are left to future work.

\paragraph{Research questions.} (RQ1) Does node/trace-level FSDS localise
injected drift more accurately than univariate PSI/KL baselines? (RQ2) Does the
covariate-versus-concept separation select the correct intervention more often
than a drift-agnostic policy? (RQ3) Does attribution-guided intervention improve
the accuracy--latency--cost frontier relative to fixed and random policies?

\paragraph{Datasets.} GSM8K and MATH (arithmetic and mathematical reasoning),
StrategyQA and HotpotQA (multi-hop question answering), HumanEval and MBPP (code
generation), and a multi-modal VQA benchmark. Drift is induced both naturally
(chronological task batches, model-version changes) and synthetically
(controlled perturbations to templates, retrieval corpora, and label mappings)
so that ground-truth drifting features are known.

\paragraph{Baselines.} CoT, SoT, ToT, and Self-Consistency as reasoning
backbones; univariate PSI/KL, fixed pruning, and random pruning as monitoring and
intervention baselines.

\paragraph{Metrics.} Task accuracy, end-to-end latency, and token cost for the
reasoning outcome; drift-detection AUC, top-$k$ recall of injected drifting
features, and Kendall's $\tau$ ranking stability for attribution quality; and
intervention lift---the change in the accuracy--latency--cost frontier---for the
end-to-end loop.

\paragraph{Splits and ablations.} Chronological splits by task batch, with nested
cross-validation (outer split by time or task, inner split for FSDS
hyperparameters). Ablations vary the outcome proxy $Y$, the nuisance
specification, the overlap threshold, and the aggregation granularity.

\section{Production Deployment}
\label{sec:deployment}

\paragraph{Batching and cadence.} FSDS runs on rolling windows of traces at a
cadence matched to traffic---every few minutes under high load, hourly or daily
otherwise---comparing the current window against a decayed reference.

\paragraph{Parallelisation.} The permutation and LOCO passes are embarrassingly
parallel and are distributed with joblib or multiprocessing; the post-fit ranking
is $O(p)$ and negligible.

\paragraph{Fallback.} Under poor overlap (below $0.3$) or a small live batch
($n<50$), the concept plane is suppressed and the system reports the covariate
plane or univariate PSI/KL, in line with Assumption~\ref{ass:overlap}.

\paragraph{Routing.} A gated mixture over reasoning experts---SoT, ToT, CoT,
direct answering, and retrieval-augmented generation---consumes the missingness
mask, quality signals, and drift shares as gate inputs, routing each request to
the policy best matched to the current regime.

\paragraph{Monitoring.} A dashboard surfaces the top-$k$ drift nodes, the
covariate and concept shares with confidence intervals, and the recommended
actions, turning the statistical procedure into an operational diagnostic.

\section{Limitations}
\label{sec:limitations}

The prototype inherits the assumptions and scope of FSDS. Attribution quality
depends on the feature map $\phi$; poorly chosen features cannot be recovered by
the estimators. The concept plane requires an outcome proxy $Y$, and the meaning
of concept drift is defined by that choice; a miscalibrated proxy biases the
concept attribution, mitigated but not eliminated by the double robustness of
\eqref{eq:porisk}. The procedure requires adequate overlap
(Assumption~\ref{ass:overlap}) and abstains otherwise. Finally, the attribution
is a distance estimate and not a causal explanation (Remark~\ref{rem:distance}),
and the tables in this manuscript specify an output format rather than reporting
measured effects.

\section{Conclusion}
\label{sec:conclusion}

We have drafted a prototype that applies Feature Selection for Distribution Shift
to Skeleton-of-Thought and Tree-of-Thoughts reasoning. By representing a
reasoning trace as a structured feature object, the procedure detects covariate
shift and concept drift at the node and trace level, ranks the features that
drive the shift, and converts the top-ranked features into concrete
interventions---re-planning, re-expansion, budget adjustment, value
re-calibration, routing, or fallback. The prototype is model-agnostic, robust to
missing branches and modalities, and designed for continuous production
monitoring, providing a principled bridge from silent reasoning drift to
actionable control.

\appendix
% Supplement: FSDS-SoT landing (decomposability, budget executor, dashboard, procedures)
% Compile standalone or \input after main manuscript macros are loaded.

\section{FSDS-SoT Landing: Procedures and Dashboard}
\label{sec:sot-landing}

We view production SoT control as five procedures. Procedures~(1)--(2) below
are the primary actionable pair discussed in the prototype; (0), (3)--(4) close
the loop.

\begin{enumerate}[leftmargin=*]
  \item[\textbf{(0)}] \textbf{Batch covariate monitoring:} $\Dold$ vs.\ $\Dnew$ on
  trace features $\phi(\mathcal{S})$; RF-domain $\VIMP$ and $\MMD$-LOCO (no $Y$).
  \item[\textbf{(1)}] \textbf{Grouping granularity:} decomposability on skeleton-point
  embeddings $\{\z_b\}_{b=1}^B$; suggested cluster count $K$ and within-cluster order.
  \item[\textbf{(2)}] \textbf{Critical-path budget executor:} per-branch expansion
  tokens $L_b$, model tier, and check budget from shift and quality.
  \item[\textbf{(3)}] \textbf{Merge / aggregation:} deduplicate and reconcile parallel
  expansions (GoT-style logic at merge time).
  \item[\textbf{(4)}] \textbf{Confidence and fallback:} overlap, bootstrap CI, stability;
  abstain or covariate-only reporting when unreliable.
\end{enumerate}

\subsection{Actionable outputs at a glance (tables)}
\label{sec:action-tables}

Table~\ref{tab:obj1-topology} summarises the primary topology object: suggested
parallel width $K$, merge rules, and execution order inside each cluster.
Table~\ref{tab:obj2-budget} gives the budget executor outputs.
Table~\ref{tab:incremental-value} lists incremental value beyond cost reduction.
Table~\ref{tab:dashboard-panels} maps the four-panel dashboard to decisions.

\begin{table}[t]
\centering
\small
\caption{Object~1 --- grouping granularity (embedding topology). $K$ is a
\emph{suggestion}; $K_{\mathrm{final}}=\mathrm{clip}(K_{\mathrm{FSDS}},K_{\min}^{\mathrm{biz}},K_{\max}^{\mathrm{biz}})$.}
\label{tab:obj1-topology}
\begin{tabular}{@{}p{0.22\linewidth}p{0.30\linewidth}p{0.38\linewidth}@{}}
\toprule
Output & How it is computed & Actionable insight \\
\midrule
Suggested size $K$ &
Coupling graph on $\{\z_b\}$; connected components / spectral clusters &
How many \emph{parallel groups} to run; clip by business min/max width \\
Cluster membership &
Cluster labels per skeleton point $b$ &
Which branches share a parallel slot vs.\ must not race \\
Merge pairs &
High coupling + concept-redundant ($C_{ij}\le\varepsilon_c$) &
Collapse duplicate skeleton points before expand (no double-counting) \\
Sequential order &
Topological / dependency order \emph{within} a cluster &
Ranked expand sequence when intra-cluster coupling is high \\
Parallel vs.\ sequential &
$P=1-\overline{H}_{ij}$, redundancy $R$, balance $\mathrm{Bal}$ &
Allow SoT parallel only if gate passes; else CoT or cluster-internal chain \\
Straggler flag &
Imbalanced predicted shift or length across $b$ &
Split or cap a branch before it dominates critical path \\
\bottomrule
\end{tabular}
\end{table}

\begin{table}[t]
\centering
\small
\caption{Object~2 --- critical-path budget executor (after topology is fixed).}
\label{tab:obj2-budget}
\begin{tabular}{@{}p{0.22\linewidth}p{0.30\linewidth}p{0.38\linewidth}@{}}
\toprule
Output & Driver & Actionable insight \\
\midrule
Expansion tokens $L_b$ &
Normalized shift $s_b$, quality $q_b$; clip to $C_{\mathrm{lat}}$ &
Set \texttt{max\_new\_tokens} per branch; longest $L_b$ sets span \\
Model tier &
High $s_b$ and risk $s_b(1-q_b)$ &
Route branch to small / medium / large model \\
Check budget &
Risk $\propto s_b(1-q_b)$; minimum $1$ if low shift but low $q_b$ &
Targeted verify / tool-validity calls (not uniform on every branch) \\
Zero-sum reclaim &
Low-shift, high-$q_b$ branches &
Reallocate saved tokens to drifting branches under fixed total budget \\
\bottomrule
\end{tabular}
\end{table}

\begin{table}[t]
\centering
\small
\caption{Incremental value of FSDS-SoT vs.\ uniform SoT (justify with A/B or
frontier: success $\times$ cost $\times$ span).}
\label{tab:incremental-value}
\begin{tabular}{@{}llp{0.42\linewidth}@{}}
\toprule
Category & Primary object & What improves (beyond lower token bill) \\
\midrule
Latency / SLA &
Object~2 &
Shorter critical path (span); fewer straggler timeouts under parallel decode \\
Reliability under drift &
Object~2 (+ $Y$) &
Higher success / tool-validity when retrieve--act branches shift \\
Observability &
Batch $\VIMP$ + Object~1 &
Localised drift tickets; faster MTTR vs.\ aggregate accuracy alone \\
Safety / compliance &
Object~2 &
Directed checks on high-risk cells; minimum guardrail on LoS$\cdot$LoQ \\
Capacity &
Object~1 + Object~2 &
Fewer bad parallel merges; more effective QPS per GPU slot \\
Strategy safety &
Object~1 &
Fallback to sequential when decomposability fails (avoid SoT tail risk) \\
Learning loop (slow) &
Batch monitoring &
Targeted retrain / AGOD on persistent drift slices \\
\bottomrule
\end{tabular}
\end{table}

\begin{table}[t]
\centering
\small
\caption{Dashboard panels: embedding-only vs.\ mixed signals.}
\label{tab:dashboard-panels}
\begin{tabular}{@{}llll@{}}
\toprule
Panel & Data source & Suggestion type & Typical action \\
\midrule
1 Coupling heatmap & Branch embedding & $K$, merge, sequentialise &
Re-plan groups before expand \\
2 PCA clusters & Branch embedding & Same partition as Panel~1 &
Visual audit of cluster assignment \\
3 Budget bars &
Embedding shift + quality &
$L_b$, tier, checks &
Deploy per-branch scheduler limits \\
4 Batch VIMP &
Trace features $\Dold$ vs.\ $\Dnew$ &
Default policy drift &
Adjust templates, caps, monitoring cadence \\
\bottomrule
\end{tabular}
\end{table}

\subsection{Suggested cluster count $K$ (not a business mandate)}
\label{sec:k-suggested}

On branch embeddings $\z_b$, form coupling $H_{ij}=\widehat{\mathrm{HSIC}}_{\mathrm{norm}}(\z_i,\z_j)$,
threshold at $\delta$ to obtain a graph, and let $K$ be the number of connected
components (or spectral clusters). The \emph{decomposability degree}
\begin{equation}
  \mathcal{D} = K/B \in (0,1]
\end{equation}
is a \emph{data-driven suggestion} for parallel width. Business constraints
$K_{\min}^{\mathrm{biz}}, K_{\max}^{\mathrm{biz}}$ (SLA, product shape, compliance)
clip the suggestion:
\begin{equation}
  K_{\mathrm{final}} = \mathrm{clip}(K_{\mathrm{FSDS}},\, K_{\min}^{\mathrm{biz}},\, K_{\max}^{\mathrm{biz}}).
\end{equation}
Redundancy rate $R$, parallelizability $P=1-\overline{H}_{ij}$, and balance
$\mathrm{Bal}$ gate whether parallel SoT is allowed at all; if not, fall back to
sequential CoT or cluster-internal ordering.

\subsection{Critical-path budget executor}
\label{sec:budget-exec}

Parallel SoT latency is governed by the span
\begin{equation}
  T_{\infty} = t_{\mathrm{skel}} + \max_{c\in\Pi}\sum_{b\in\mathrm{chain}(c)} L_b\, t_{\mathrm{tok}} + t_{\mathrm{merge}},
\end{equation}
not by total token count. Given normalized shift $s_b$ (from embedding distance
to skeleton reference) and branch quality $q_b$ (verifier/PRM, external to pure
embedding geometry), allocate
\begin{equation}
  L_b = \mathrm{clip}(\kappa\, s_b\, q_b,\, L_{\min},\, C_{\mathrm{lat}}),
  \qquad
  \mathrm{risk}_b \propto s_b(1-q_b),
\end{equation}
with check budget $\lceil \beta_0 + \beta_1\,\mathrm{risk}_b\rceil$ clipped to
$[1,5]$, and model tier (small/medium/large) from joint high shift and high risk.
A \emph{minimum check} of $1$ is retained when quality is low even if shift is
low (LoS$\cdot$LoQ): attribution reads ``no new content, but not trusted''---do
not upscale the model; spend a cheap guardrail check instead.

\subsection{Topology $\to$ resources visualization}
\label{sec:viz-logic}

The operational dashboard separates embedding-only topology from resource actions:
\begin{itemize}[leftmargin=*]
  \item \textbf{Panel 1 (coupling heatmap):} embedding-only; actionable insight =
  suggested $K$, merge pairs, sequentialize high-$H_{ij}$ blocks.
  \item \textbf{Panel 2 (PCA-2, cluster colors):} same partition as Panel~1.
  \item \textbf{Panel 3 (budget executor):} shift (blue) and quality (green) on
  $[0,1]$; risk $= s(1-q)$ (red); expansion tokens $L_b$ (purple, right axis);
  tier$\cdot$checks annotation. Quadrants HiS$\cdot$HiQ, HiS$\cdot$LoQ,
  LoS$\cdot$HiQ, LoS$\cdot$LoQ guide interpretation.
  \item \textbf{Panel 4 (batch $\VIMP$):} trace-level covariate drift across
  $\Dold$ vs.\ $\Dnew$; actionable insight = re-plan defaults and monitoring
  alerts, not per-request decode limits.
\end{itemize}
Embedding geometry drives Panels~1--2 and the shift component of Panel~3; quality
is typically non-embedding; Panel~4 is batch monitoring.

\subsection{ToT width and depth vs.\ SoT branches}
\label{sec:tot-vs-sot}

For ToT, \emph{width} is branching factor per node and \emph{depth} is search
horizon; FSDS adjusts where to expand or prune. For SoT, width is the suggested
$K$ and depth is per-branch $L_b$; the same budget executor applies after
topology is fixed. Tree search and branch count both answer ``how many groups,''
but ToT additionally requires value-head re-calibration when $\CD$ concentrates
on node value features.

\subsection{Intervention stack (beyond structure and content)}
\label{sec:intervention-stack}

The separation principle ($\CS$ $\to$ structure; $\CD$ $\to$ content or scoring)
extends to a four-layer stack ordered by latency-to-effect: (i)~inference-time
(budget reallocation, routing, retrieval, checks, cache); (ii)~configuration
(prompt/template, rollback); (iii)~learning/data (AGOD, targeted retraining);
(iv)~human/monitoring. Budget reallocation uses water-filling at fixed total
budget (zero-sum across branches) or raised budget under elevated batch risk.

\subsection{Economic justification (sketch)}
\label{sec:econ}

Uniform SoT sets $L_b=L_{\mathrm{uniform}}$ and a single model tier, so critical-path
latency and cost scale with the longest branch. FSDS-SoT shortens low-shift
branches, routes large tiers only to high shift$\times$risk cells, and targets
checks, improving span (p99 under parallel decoding), tokens per request, and
net savings relative to ms-level FSDS overhead (implementation:
\texttt{fsds\_sot/economics.py}).


\begin{thebibliography}{99}
\small

\bibitem{wei2022cot} J.~Wei, X.~Wang, D.~Schuurmans, M.~Bosma, B.~Ichter,
F.~Xia, E.~Chi, Q.~Le, and D.~Zhou. Chain-of-Thought Prompting Elicits Reasoning
in Large Language Models. \emph{NeurIPS}, 2022.

\bibitem{ning2024sot} X.~Ning, Z.~Lin, Z.~Zhou, Z.~Wang, H.~Yang, and Y.~Wang.
Skeleton-of-Thought: Large Language Models Can Do Parallel Decoding.
\emph{ICLR}, 2024.

\bibitem{yao2023tot} S.~Yao, D.~Yu, J.~Zhao, I.~Shafran, T.~L. Griffiths,
Y.~Cao, and K.~Narasimhan. Tree of Thoughts: Deliberate Problem Solving with
Large Language Models. \emph{NeurIPS}, 2023.

\bibitem{besta2024got} M.~Besta et al. Graph of Thoughts: Solving Elaborate
Problems with Large Language Models. \emph{AAAI}, 2024.

\bibitem{yao2023react} S.~Yao, J.~Zhao, D.~Yu, N.~Du, I.~Shafran,
K.~Narasimhan, and Y.~Cao. ReAct: Synergizing Reasoning and Acting in Language
Models. \emph{ICLR}, 2023.

\bibitem{wang2023selfconsistency} X.~Wang, J.~Wei, D.~Schuurmans, Q.~Le,
E.~Chi, S.~Narang, A.~Chowdhery, and D.~Zhou. Self-Consistency Improves Chain of
Thought Reasoning in Language Models. \emph{ICLR}, 2023.

\bibitem{shinn2023reflexion} N.~Shinn, F.~Cassano, E.~Berman, A.~Gopinath,
K.~Narasimhan, and S.~Yao. Reflexion: Language Agents with Verbal Reinforcement
Learning. \emph{NeurIPS}, 2023.

\bibitem{gretton2012mmd} A.~Gretton, K.~Borgwardt, M.~Rasch, B.~Sch\"olkopf,
and A.~Smola. A Kernel Two-Sample Test. \emph{JMLR}, 13:723--773, 2012.

\bibitem{gretton2005hsic} A.~Gretton, O.~Bousquet, A.~Smola, and
B.~Sch\"olkopf. Measuring Statistical Dependence with Hilbert--Schmidt Norms.
\emph{ALT}, 2005.

\bibitem{candes2018knockoffs} E.~Cand\`es, Y.~Fan, L.~Janson, and J.~Lv.
Panning for Gold: Model-X Knockoffs for High-Dimensional Controlled Variable
Selection. \emph{JRSS-B}, 80(3):551--577, 2018.

\bibitem{nie2021rlearner} X.~Nie and S.~Wager. Quasi-Oracle Estimation of
Heterogeneous Treatment Effects. \emph{Biometrika}, 108(2):299--319, 2021.

\bibitem{kennedy2023drlearner} E.~H. Kennedy. Towards Optimal Doubly Robust
Estimation of Heterogeneous Causal Effects. \emph{Electronic Journal of
Statistics}, 2023.

\bibitem{wager2018causalforest} S.~Wager and S.~Athey. Estimation and Inference
of Heterogeneous Treatment Effects using Random Forests. \emph{JASA},
113(523):1228--1242, 2018.

\bibitem{breiman2001rf} L.~Breiman. Random Forests. \emph{Machine Learning},
45(1):5--32, 2001.

\bibitem{fsds2026} Feature Selection for Distribution Shift: Meta-Learner
Objective Functions via Pseudo-Outcome and R-Learners. Companion manuscript,
2026.

\end{thebibliography}

\end{document}
